% GENERATED FILE. Edit canonical lessons, then run npm run generate:books.
% canonical-source-hash: fnv1a64:ced4ff3eb6183388
% canonical-lessons: ML-C361-yugam, ML-C361-bhutakalam, ML-C361-sandarbham, ML-C361-tutakkam, ML-C361-apratiksitamayi

\chapter{Ages and Occasions}
\label{ch:ml-ages-and-occasions}
\hlchaptermodality{361}

\begin{chapteropening}
\textbf{By the end of this chapter:} \emph{I can say an era, the past, an occasion, a beginning and unexpectedly.}

Complete the last lesson of chapter 361: I can say an era, the past, an occasion, a beginning and unexpectedly.
\end{chapteropening}

\section[\texorpdfstring{\ml{യുഗം} (yugaṁ)}{യുഗം (yugaṁ)}]{\ml{യുഗം} (yugaṁ) — an era; an age}

\label{lesson:ML-C361-yugam}



\begin{quote}
Before the new one: say the Malayalam for luck, then the Malayalam for as fate willed.
\end{quote}

\subsection*{You'll want to know: \ml{യുഗം}}
\textbf{\ml{യുഗം}} — \emph{yugaṁ} — \textquotedblleft{}an era; an age\textquotedblright{}.

\begin{quote}
\ml{ഇത്} \ml{ഇന്റർനെറ്റിന്റെ} \ml{യുഗമാണ്}.
\end{quote}

\begin{quote}
\emph{itŭ inṟarneṟṟinṟe yugamāṇŭ.}
\end{quote}

\begin{quote}
This is the age of the internet.
\end{quote}

In old tales a \ml{യുഗം} is one of the great ages of the world; today it is any era.

\begin{cousinweb}
\textbf{\ml{യുഗം}} (\emph{yugaṁ}) is borrowed from Sanskrit \emph{yuga}, first \textquotedblleft{}a yoke, a pair\textquotedblright{} (CDIAL 10482), and in later Sanskrit also an age of the world (Monier-Williams).
\end{cousinweb}

\begin{grammarlens}[title={What you've built}]
One more word for telling what happened.
\end{grammarlens}

\subsection*{Your turn}
\begin{itemize}
\raggedright
  \item \emph{Say it:} \emph{yugaṁ}
  \item \emph{Say it:} \emph{yugaṁ}, once more
  \item \emph{Say it:} \emph{itŭ inṟarneṟṟinṟe yugamāṇŭ.} — This is the age of the internet.
  \item \emph{Recall:} say the Malayalam for luck, then the Malayalam for as fate willed, then say \emph{yugaṁ} again
\end{itemize}

\begin{tcolorbox}[breakable,colback=teal!4,colframe=teal!35!black,title={Before you move on}]
What does \ml{യുഗം} mean? (\textquotedblleft{}An era; an age\textquotedblright{}.) Say it once more.
\end{tcolorbox}
\section[\texorpdfstring{\ml{ഭൂതകാലം} (bhūtakālaṁ)}{ഭൂതകാലം (bhūtakālaṁ)}]{\ml{ഭൂതകാലം} (bhūtakālaṁ) — the past}

\label{lesson:ML-C361-bhutakalam}



\begin{quote}
Before the new one: say the Malayalam for as fate willed, then the Malayalam for an era.
\end{quote}

\subsection*{You'll want to know: \ml{ഭൂതകാലം}}
\textbf{\ml{ഭൂതകാലം}} — \emph{bhūtakālaṁ} — \textquotedblleft{}the past\textquotedblright{}.

\begin{quote}
\ml{അവൻ} \ml{ഭൂതകാലത്തെക്കുറിച്ച്} \ml{ഒന്നും} \ml{പറഞ്ഞില്ല}.
\end{quote}

\begin{quote}
\emph{avan bhūtakālattekkuṟiccŭ onnuṁ paṟaññilla.}
\end{quote}

\begin{quote}
He said nothing about the past.
\end{quote}

In grammar, the past tense. \ml{ഭാവി}, the future, is its opposite.

\begin{cousinweb}
\textbf{\ml{ഭൂതകാലം}} (\emph{bhūtakālaṁ}) is borrowed from Sanskrit \emph{bhūtakāla}, \textquotedblleft{}past time; the past tense\textquotedblright{} (Monier-Williams): \emph{bhūta}, \textquotedblleft{}what has been\textquotedblright{}, and \emph{kāla}, \textquotedblleft{}time\textquotedblright{} (CDIAL 3084), the word that gave \ml{കാലം}.
\end{cousinweb}

\begin{grammarlens}[title={What you've built}]
One more word for telling what happened.
\end{grammarlens}

\subsection*{Your turn}
\begin{itemize}
\raggedright
  \item \emph{Say it:} \emph{bhūtakālaṁ}
  \item \emph{Say it:} \emph{bhūtakālaṁ}, once more
  \item \emph{Say it:} \emph{avan bhūtakālattekkuṟiccŭ onnuṁ paṟaññilla.} — He said nothing about the past.
  \item \emph{Recall:} say the Malayalam for as fate willed, then the Malayalam for an era, then say \emph{bhūtakālaṁ} again
\end{itemize}

\begin{tcolorbox}[breakable,colback=teal!4,colframe=teal!35!black,title={Before you move on}]
What does \ml{ഭൂതകാലം} mean? (\textquotedblleft{}The past\textquotedblright{}.) Say it once more.
\end{tcolorbox}
\section[\texorpdfstring{\ml{സന്ദർഭം} (sandarbhaṁ)}{സന്ദർഭം (sandarbhaṁ)}]{\ml{സന്ദർഭം} (sandarbhaṁ) — an occasion; a context}

\label{lesson:ML-C361-sandarbham}



\begin{quote}
Before the new one: say the Malayalam for an era, then the Malayalam for the past.
\end{quote}

\subsection*{You'll want to know: \ml{സന്ദർഭം}}
\textbf{\ml{സന്ദർഭം}} — \emph{sandarbhaṁ} — \textquotedblleft{}an occasion; a context\textquotedblright{}.

\begin{quote}
\ml{ആ} \ml{സന്ദർഭത്തിൽ} \ml{ഞാൻ} \ml{ഒന്നും} \ml{പറഞ്ഞില്ല}.
\end{quote}

\begin{quote}
\emph{ā sandarbhattil ñān onnuṁ paṟaññilla.}
\end{quote}

\begin{quote}
On that occasion I said nothing.
\end{quote}

Also the context of a word or a sentence.

\begin{cousinweb}
\textbf{\ml{സന്ദർഭം}} (\emph{sandarbhaṁ}) is borrowed from Sanskrit \emph{saṁdarbha}, \textquotedblleft{}a stringing or weaving together; a composition\textquotedblright{} (Monier-Williams). An occasion is how things are strung together at one time.
\end{cousinweb}

\begin{grammarlens}[title={What you've built}]
One more word for telling what happened.
\end{grammarlens}

\subsection*{Your turn}
\begin{itemize}
\raggedright
  \item \emph{Say it:} \emph{sandarbhaṁ}
  \item \emph{Say it:} \emph{sandarbhaṁ}, once more
  \item \emph{Say it:} \emph{ā sandarbhattil ñān onnuṁ paṟaññilla.} — On that occasion I said nothing.
  \item \emph{Recall:} say the Malayalam for an era, then the Malayalam for the past, then say \emph{sandarbhaṁ} again
\end{itemize}

\begin{tcolorbox}[breakable,colback=teal!4,colframe=teal!35!black,title={Before you move on}]
What does \ml{സന്ദർഭം} mean? (\textquotedblleft{}An occasion; a context\textquotedblright{}.) Say it once more.
\end{tcolorbox}
\section[\texorpdfstring{\ml{തുടക്കം} (tuṭakkaṁ)}{തുടക്കം (tuṭakkaṁ)}]{\ml{തുടക്കം} (tuṭakkaṁ) — a beginning; a start}

\label{lesson:ML-C361-tutakkam}



\begin{quote}
Before the new one: say the Malayalam for the past, then the Malayalam for an occasion.
\end{quote}

\subsection*{You'll want to know: \ml{തുടക്കം}}
\textbf{\ml{തുടക്കം}} — \emph{tuṭakkaṁ} — \textquotedblleft{}a beginning; a start\textquotedblright{}.

\begin{quote}
\ml{കഥയുടെ} \ml{തുടക്കം} \ml{വളരെ} \ml{രസകരമായിരുന്നു}.
\end{quote}

\begin{quote}
\emph{kathayuṭe tuṭakkaṁ vaḷare rasakaramāyirunnu.}
\end{quote}

\begin{quote}
The beginning of the story was very interesting.
\end{quote}

The noun of \ml{തുടങ്ങുക}, to begin; its opposite is \ml{ഒടുക്കം}, the end.

\begin{cousinweb}
\textbf{\ml{തുടക്കം}} (\emph{tuṭakkaṁ}) is native Dravidian, the noun of \ml{തുടങ്ങുക} (\emph{tuṭaṅṅuka}), \textquotedblleft{}to begin\textquotedblright{} (DEDR 3481), as \ml{ഒടുക്കം} is the noun of \ml{ഒടുങ്ങുക}.
\end{cousinweb}

\begin{grammarlens}[title={What you've built}]
One more word for telling what happened.
\end{grammarlens}

\subsection*{Your turn}
\begin{itemize}
\raggedright
  \item \emph{Say it:} \emph{tuṭakkaṁ}
  \item \emph{Say it:} \emph{tuṭakkaṁ}, once more
  \item \emph{Say it:} \emph{kathayuṭe tuṭakkaṁ vaḷare rasakaramāyirunnu.} — The beginning of the story was very interesting.
  \item \emph{Recall:} say the Malayalam for the past, then the Malayalam for an occasion, then say \emph{tuṭakkaṁ} again
\end{itemize}

\begin{tcolorbox}[breakable,colback=teal!4,colframe=teal!35!black,title={Before you move on}]
What does \ml{തുടക്കം} mean? (\textquotedblleft{}A beginning; a start\textquotedblright{}.) Say it once more.
\end{tcolorbox}
\section[\texorpdfstring{\ml{അപ്രതീക്ഷിതമായി} (apratīkṣitamāyi)}{അപ്രതീക്ഷിതമായി (apratīkṣitamāyi)}]{\ml{അപ്രതീക്ഷിതമായി} (apratīkṣitamāyi) — unexpectedly}

\label{lesson:ML-C361-apratiksitamayi}



\begin{quote}
Before the new one: say the Malayalam for an occasion, then the Malayalam for a beginning.
\end{quote}

\subsection*{You'll want to know: \ml{അപ്രതീക്ഷിതമായി}}
\textbf{\ml{അപ്രതീക്ഷിതമായി}} — \emph{apratīkṣitamāyi} — \textquotedblleft{}unexpectedly\textquotedblright{}.

\begin{quote}
\ml{അമ്മാവൻ} \ml{ഇന്നലെ} \ml{അപ്രതീക്ഷിതമായി} \ml{വീട്ടിൽ} \ml{വന്നു}.
\end{quote}

\begin{quote}
\emph{ammāvan innale apratīkṣitamāyi vīṭṭil vannu.}
\end{quote}

\begin{quote}
Uncle came home unexpectedly yesterday.
\end{quote}

From \ml{പ്രതീക്ഷിക്കുക}, to expect: \ml{അ}-, not, + expected + -\ml{ആയി}.

\begin{cousinweb}
\textbf{\ml{അപ്രതീക്ഷിതമായി}} (\emph{apratīkṣitamāyi}) is built from Sanskrit \emph{pratīkṣita}, \textquotedblleft{}expected\textquotedblright{} (CDIAL 8621), with \emph{a-}, \textquotedblleft{}not\textquotedblright{}, in front and the Malayalam adverb ending -\ml{ആയി}. It is a cousin of \ml{പ്രതീക്ഷ} (\emph{pratīkṣa}), \textquotedblleft{}hope\textquotedblright{}.
\end{cousinweb}

\begin{grammarlens}[title={What you've built}]
One more word for telling what happened.
\end{grammarlens}

\subsection*{Your turn}
\begin{itemize}
\raggedright
  \item \emph{Say it:} \emph{apratīkṣitamāyi}
  \item \emph{Say it:} \emph{apratīkṣitamāyi}, once more
  \item \emph{Say it:} \emph{ammāvan innale apratīkṣitamāyi vīṭṭil vannu.} — Uncle came home unexpectedly yesterday.
  \item \emph{Recall:} say the Malayalam for an occasion, then the Malayalam for a beginning, then say \emph{apratīkṣitamāyi} again
\end{itemize}

\begin{tcolorbox}[breakable,colback=teal!4,colframe=teal!35!black,title={Before you move on}]
What does \ml{അപ്രതീക്ഷിതമായി} mean? (\textquotedblleft{}Unexpectedly\textquotedblright{}.) Say it once more.
\end{tcolorbox}
